\documentclass[11pt]{article}
\usepackage[a4paper,margin=1in]{geometry}
\usepackage{amsmath,amssymb,amsthm,mathtools}
\usepackage{booktabs}
\usepackage{graphicx}
\usepackage{hyperref}
\usepackage{cleveref}

\title{Tencent (0700.HK) V4 Valuation Snapshot:\\A Multi-Method Investment Note}
\author{Tencent Model Project}
\date{As-of 2026-04-02}

\newcommand{\E}{\mathbb{E}}
\newcommand{\Var}{\mathrm{Var}}
\DeclareMathOperator*{\argmax}{arg\,max}

\begin{document}
\maketitle

\begin{abstract}
This note summarizes the V4 valuation stack for Tencent and reports a publication-ready
snapshot with code-verified outputs, quality gates, and a visual evidence pack.
At the 2026-04-02 snapshot, the model indicates market price exceeds base and
expected intrinsic values under the maintained assumptions.
\end{abstract}

\section{Model Setup}

\subsection{Discounted Cash Flow}
For scenario $s$, enterprise value is:
\begin{equation}
EV_s = \sum_{t=1}^{T} \frac{FCFF_{s,t}}{(1+WACC)^t}
+ \frac{TV_s}{(1+WACC)^T},
\end{equation}
with Gordon terminal value
\begin{equation}
TV_s = \frac{FCFF_{s,T+1}}{WACC - g_s}, \qquad FCFF_{s,T+1}=FCFF_{s,T}(1+g_s).
\end{equation}
Equity value is
\begin{equation}
EqV_s = EV_s - Debt + Cash - MinorityInterest,
\end{equation}
and fair value per share is $P^{\ast}_s = EqV_s/N_{\text{shares}}$.

\subsection{Cost of Capital}
Official cost of equity uses CAPM with country risk premium:
\begin{equation}
R_e^{\mathrm{CAPM}} = R_f + \beta_L \cdot ERP + CRP.
\end{equation}
WACC is:
\begin{equation}
WACC = \frac{E}{D+E}R_e + \frac{D}{D+E}R_d(1-\tau).
\end{equation}
APT estimates are included as diagnostics and excluded from headline discount rates when
stability gates fail.

\subsection{Ensemble}
Headline expected fair value aggregates methods $m \in \mathcal{M}$:
\begin{equation}
P^{\ast}_{\text{ens},s} = \sum_{m \in \mathcal{M}} \omega_m P^{\ast}_{m,s},
\qquad \sum_{m \in \mathcal{M}} \omega_m = 1,\quad \omega_m \ge 0.
\end{equation}

\section{Snapshot Results}

Market price is HKD 489.20/share. Base DCF fair value is HKD 354.36/share and expected
ensemble fair value is HKD 295.21/share, both implying negative margin of safety.

\begin{table}[h]
\centering
\begin{tabular}{lrrr}
\toprule
Scenario & DCF (HKD/share) & Ensemble (HKD/share) & MoS vs Market \\
\midrule
Base & 354.36 & 361.08 & -27.6\% / -26.2\% \\
Bad & 206.34 & 244.64 & -57.8\% / -50.0\% \\
Extreme & 147.87 & 193.66 & -69.8\% / -60.4\% \\
\bottomrule
\end{tabular}
\caption{DCF and ensemble valuation outputs for the 2026-04-02 snapshot.}
\end{table}

\section{Quality and Identification Caveats}

The QA summary reports 27 checks, 2 warnings, and 2 failures; investor-grade status is false.
Main caveats are CAPM--APT divergence, elevated ERP assumptions, and weak backtest
calibration strength from limited sample depth.

\section{Visual Evidence}

\begin{figure}[h]
\centering
\includegraphics[width=0.84\textwidth]{../figures/2026-04-02/01_dcf_vs_market.png}
\caption{DCF fair value by scenario versus market price.}
\end{figure}

\begin{figure}[h]
\centering
\includegraphics[width=0.84\textwidth]{../figures/2026-04-02/05_monte_carlo_distribution.png}
\caption{Monte Carlo fair-value distribution with market overlay.}
\end{figure}

\begin{figure}[h]
\centering
\includegraphics[width=0.84\textwidth]{../figures/2026-04-02/09_backtest_scatter.png}
\caption{Backtest signal-strength vs forward 12M return.}
\end{figure}

\section{Conclusion}

Under current V4 assumptions and diagnostics, the 2026-04-02 snapshot does not support
a margin-of-safety buy at market price. A bullish conclusion requires both stronger
fundamental assumptions and improved empirical calibration quality.

\end{document}
